\chapter{从局部估计到全局一致：SLAM系统}

\label{chap:md-12}

一台机器人从房间出发，沿走廊行驶，绕过一圈后回到房间。IMU不断报告短时间运动，LiDAR逐帧看到墙、地面和门框。仅靠相邻时刻的局部估计，机器人能持续给出“我刚才怎样动了”；但小误差会逐渐积累，走回房间时，估计的终点可能没有落回起点。完整SLAM系统必须同时做三件事：及时给出当前位姿，保留历史测量以便重新调整，并识别再次到达的旧地点。

第\ref{chap:md-08}章与第\ref{chap:md-09}章讨论当前状态怎样随数据更新；第\ref{chap:md-10}章保存关键帧之间的惯性相对测量；第\ref{chap:md-11}章从点云获得几何相对位姿。本章把它们放进同一台LiDAR–IMU机器人的时间线上。先分清系统保存的对象，再依次追踪前端、后端、回环和全局校正返回前端的过程。

\section{同一台机器人里有哪些长期存在的对象}

\subsection{高频状态、关键帧状态和地图}

实时估计器在每个IMU时刻维护当前状态$(\hat x_t,P_t)$，用于去畸变、控制或显示。后端选择较少的时刻$t_i$作为关键帧，保存$x_i=(R_i,p_i,v_i,b_{g,i},b_{a,i})$或其中与任务相关的变量。地图则由点云、局部表面或关键帧子地图构成。当前状态可以高频变化，关键帧集合只在必要时增长，地图还要记录点云由哪些位姿放入同一坐标系。

\begin{quote}\textbf{同一个位姿的三种用途}\par

当前位姿是实时输出；关键帧位姿是后端要反复调整的变量；用于拼接地图的位姿是点云坐标变换的依据。三者可以在某个时刻取相同数值，但保存时不能只留下一个没有来源的矩阵。后端修改关键帧后，地图和当前输出如何跟随，也需要有明确规则。

\end{quote}

\subsection{观测、测量摘要与因子}

原始IMU读数和原始点云是传感器数据。预积分记录和GICP配准结果是由数据计算出的摘要。因子则进一步规定：给定哪些待估状态，这份摘要应产生什么残差和权重。例如预积分记录可以反复评价$(x_i,x_j)$；GICP得到的相对位姿可以比较$T_i^{-1}T_j$。一份摘要也可能只用于给另一个算法提供初值，而不成为新的独立因子。

\subsection{三个模块各回答什么问题}

前端需要及时回答“从刚才到现在怎样动了”；后端要找一组能共同解释历史测量的关键帧状态；回环要检索并验证“当前看到的是不是以前来过的地方”。回环通过几何验证产生远距离约束，最终仍交由后端与其他因子一起处理。

\section{前端：从原始数据得到当前运动}

\subsection{时间戳、外参与坐标链先于融合}

设$W$是地图所用世界系，$B$是机体系，$L$是LiDAR系。IMU动力学、点云坐标和关键帧位姿之间必须有固定的外参变换，如$T_{BL}$，并知道每个点及每个IMU样本的采集时刻。若时间戳相差一段固定延迟，快速转动时点云会按错误姿态投影；这类误差未必能靠后端调大噪声解决。外参若未知或待估，应作为明确参数进入模型，不能在不同模块中各自假定一个值。

\subsection{IMU预测怎样提供扫描期间的运动}

扫描期间IMU频繁到来，ESKF按第八章的名义状态递推和协方差递推产生连续的短时预测。它为下一帧点云提供位姿初值，也为扫描内不同采样时刻提供运动估计。这里只把IMU预测看成当前可用的参考；偏置误差和噪声仍会使预测随时间偏离。

\subsection{点云去畸变为什么在配准之前}

一帧LiDAR点云通常不是在同一瞬间获得。点$p_m$的时间戳为$t_m$，若要把整帧都表达在选定的参考时刻$t_0$，需要用扫描期间的估计轨迹将$p_m$从当时传感器系变到参考系。机器人在扫描时转动较快，直接把所有点当成$t_0$采集，会使本来平直的墙产生弯曲或重影。GICP随后可能试图用一个刚体位姿解释这种非刚体形变，导致错误配准。

\subsection{GICP怎样修正局部几何位姿}

去畸变后，以IMU预测或上一帧里程计作为GICP初值，将当前点云与上一关键帧或局部子地图配准。前端检查对应数量、重叠率、残差和退化方向，再决定是否接受结果。配准返回的相对位姿可作为ESKF的一次外部观测；另一种设计直接用点到局部表面的残差迭代更新状态。这两种接口要选定一种主线，且要让观测噪声与配准质量对应。

\subsection{何时保存关键帧}

每帧都加入后端会使状态数和点云数据不断膨胀；只按固定时间保存又可能错过新几何。可综合平移、转角、经过的时间、与现有子地图的重叠率及是否看到新的结构。关键帧保存时间戳、状态初值、用于地图的点云或局部摘要、必要的原始IMU区间和回环检索描述子。关键帧选择改变后端变量集合，也决定预积分区间的端点。

\section{前端交给后端的究竟是什么}

\subsection{IMU区间记录不是当前状态协方差}

相邻关键帧$i,j$之间，第十章从零开始累积$\widetilde{\Delta R}_{ij},\widetilde{\Delta v}_{ij},\widetilde{\Delta p}_{ij}$、协方差和偏置Jacobian。它们描述这段IMU测量的误差；ESKF当前状态协方差$P_t$则已经包含此前估计和更新的信息。若把两者当成同一份独立测量加入后端，会重复计算信息。

\subsection{GICP输出是初始化还是测量}

如果GICP只用来给关键帧位姿提供初值，后端不必再为这个结果建立因子。若它作为几何相对位姿测量，就需要明确源/目标方向、残差局部坐标、有效协方差和配准质量。若同一组点已经直接在后端构成点云因子，再把由这些点压缩出的GICP位姿当作独立因子，可能重复使用几何信息。

\subsection{使用IMU初值不自动产生独立性}

GICP可以用IMU预测决定初始对应；这提高配准成功率，但配准输出的误差模型未必与IMU因子严格独立。是否近似为条件独立，需要说明条件化了哪些数据、配准是否只用初值、失败样本怎样筛选。系统实现可采用近似，但不能把“有两个程序输出”直接等同于“两份独立证据”。

\section{后端：多份局部约束怎样形成联合问题}

\subsection{节点与因子的最小例子}

取三个关键帧$x_0,x_1,x_2$。预积分因子连接相邻帧$(x_0,x_1)$与$(x_1,x_2)$；几何因子可以连接配准过的两帧；一个先验因子固定初始参考。每个因子只接触少数变量，因此整体Jacobian是稀疏的。若以后在$x_0,x_2$之间发现回环，就新增一条跨越较长时间的几何边。

\subsection{先从概率模型得到加权残差}

在选定的局部高斯和条件独立近似下，后验与各因子的乘积成正比。取负对数后，待优化目标可写为

\begin{equation*}
E(X)=E_{\mathrm{prior}}(X)
+\sum_{(i,j)\in\mathcal E_I}\rho_I\!\left(\|r_{ij}^{\mathrm{IMU}}\|_{\Sigma_{ij}^{-1}}^2\right)
+\sum_{(i,j)\in\mathcal E_G}\rho_G\!\left(\|r_{ij}^{\mathrm{geom}}\|_{\Lambda_{ij}^{-1}}^2\right)
+\sum_{(i,j)\in\mathcal E_L}\rho_L\!\left(\|r_{ij}^{\mathrm{loop}}\|_{\Omega_{ij}^{-1}}^2\right).
\end{equation*}

这里$X$是全部待估关键帧变量，$\mathcal E_I,\mathcal E_G,\mathcal E_L$分别表示惯性、普通几何和回环约束。协方差的逆只在残差坐标、顺序和单位都匹配时才是正确权重。鲁棒核在相应因子可能有离群值时使用，具体形式也属于模型选择。

\subsection{为什么必须给整体位姿一个参考}

如果所有观测只比较相对运动，把整条轨迹与整张地图同时作同一个全局刚体变换，相对残差不变。于是全局原点和朝向的若干方向无法仅由这些观测确定，线性化Hessian也会出现零方向。常见处理是固定第一帧位姿或加入明确的位姿先验。重力、GNSS和其他绝对信息会改变具体的可观性；不能不看传感器组合就说“必然恰好有六个自由方向”。

\subsection{一次局部优化怎样更新流形状态}

在当前$X$附近给每个姿态选择合法局部扰动，得到$r_a(X\oplus\delta X)\simeq r_a(X)+J_a\delta X$。白化后汇总各因子，形成稀疏的局部二次模型并求$\delta X$。随后使用指数映射更新姿态，位置、速度和偏置按所选欧氏坐标更新，再重新计算残差与Jacobian。第\ref{chap:md-06}章说明了为何必须这样在局部线性空间中算增量，再回到非线性状态空间。

\subsection{滑动窗口与边缘化}

为了使实时计算量可控，局部后端可只优化最近一段关键帧。移出窗口的旧状态对保留状态的影响，可以在线性化点附近经Schur补形成先验。这个先验保存了被移除变量带来的局部信息，也继承了当时的线性化约定；不能把它理解为与线性化点无关、可随意搬动的精确原始测量。全局位姿图则可以在较低频率调整更长的历史轨迹。

\section{回环：从“像以前的地方”到可信约束}

\subsection{地点识别先找候选}

机器人来到一处走廊入口，当前点云与历史某个子地图可能相似。可以为关键帧保存便于检索的场景描述，再在历史库中找候选；\href{https://arxiv.org/abs/1805.03465}{Scan Context论文}提供了一种LiDAR地点识别方法。检索阶段可排除时间上过近的关键帧，因为它们通常已由普通里程计连接。描述子相似只说明值得检查，不能直接作为六维位姿因子。

\subsection{几何验证怎样接上GICP}

对候选历史子地图给出粗初值，用GICP做精配准，并检查有效重叠、残差、稳定性及弱约束方向。必要时还可以要求连续多个关键帧在同一历史区域得到相容结果。只有通过验证，才保存测得的回环相对变换$Z_{ij}^{\mathrm{loop}}$与相应权重。重复结构环境中，几何局部一致仍可能对应错误地点，因此验证应结合更广的时空信息。

\subsection{回环因子怎样拉动历史状态}

若当前优化预测的相对位姿为$T_i^{-1}T_j$，回环残差可写成

\begin{equation*}
r_{ij}^{\mathrm{loop}}
=\operatorname{Log}\!\left((Z_{ij}^{\mathrm{loop}})^{-1}T_i^{-1}T_j\right)^\vee.
\end{equation*}

当轨迹累积漂移时，这条边往往与原来的相邻边不能同时完全满足。后端根据各因子的协方差和鲁棒代价，在多帧之间分配修正，而不是直接把$j$帧“设成”$i$帧。错误回环能破坏整张地图，因此可使用鲁棒核、可切换约束或事前更严格的几何筛选。

\section{全局修正怎样返回实时系统}

\subsection{关键帧改了，地图也要跟着改}

全局优化改变历史关键帧位姿后，由这些位姿拼出的点云地图也需重新投影或更新子地图位姿。若只修改轨迹数组，地图点仍停在旧坐标，显示与后续配准将互相矛盾。对大型地图，可保留局部子地图并更新它们的锚定位姿，不必每次重建全部原始点。

\subsection{连续里程计坐标与可校正地图坐标}

前端通常需要连续的高频输出，回环却可能使全局地图位姿突然改变。可以区分局部连续的里程计坐标系$O$和全局可校正的地图坐标系$W$：前端持续给出$T_{OB}(t)$，后端更新$T_{WO}$或等价的全局对齐关系，供地图与全局查询使用。这里的分离是一种系统组织方式；具体实现还需明确$T_{WO}$随时间如何更新、历史输出如何解释。

\subsection{后端状态怎样成为下一轮前端的参考}

前端可在适当时刻使用后端优化后的关键帧或偏置估计修正自身参考，但必须处理时间延迟和已经传播出去的IMU数据。若直接用旧关键帧结果覆盖当前高频状态，可能产生时间不连续。更稳妥的做法是明确状态重定位、重传播或坐标变换规则，并保证后续点云去畸变使用与当前状态一致的轨迹。

\section{沿时间走完一次系统流程}

\subsection{从第一帧到下一关键帧}

首先确定世界参考、重力、初始姿态和偏置估计；若这些量尚未可靠初始化，后面的因子权重与去畸变都可能出错。每个IMU样本推动ESKF当前状态，LiDAR点按采集时刻去畸变。当前扫描在IMU初值附近与局部地图配准，几何结果经过质量检查后修正当前估计。运动量或地图覆盖变化达到条件时，保存关键帧。

\subsection{从关键帧到局部后端}

对相邻关键帧区间预积分，得到相对量、协方差和偏置Jacobian；保存合格的几何相对位姿或原始点云因子。局部后端以当前关键帧状态为初值，加入先验及相邻因子，反复评价端点残差并优化。偏置稍变时按第十章一阶修正预积分，变化大时重新积分。优化后的关键帧和局部子地图结果返回前端作为后续参考。

\subsection{从再次到达旧地点到全局一致}

地点识别从历史关键帧中找到候选，GICP在相应子地图间做几何验证。若通过检查，加入远距离回环因子，运行全局优化。历史关键帧随之移动，地图重新对齐，实时状态通过既定坐标关系接到新地图。此时系统没有改写过去的传感器读数；改变的是对整条轨迹和地图的联合估计。

\section{几种常见失效在数学上表现为什么}

\begin{center}\small\begin{tabular}{p{0.420\textwidth}p{0.420\textwidth}}\toprule

现象 & 需要检查的对象 \\

\midrule

转弯时墙面被拉弯 & 时间同步、扫描内运动估计和外参 \\

配准残差小却沿走廊持续漂移 & 局部Hessian的弱方向、子地图几何与权重 \\

后端比前端还“确信”错误位姿 & 协方差过小、错误对应及同一数据重复计数 \\

回环后地图断裂或重复 & 错误地点候选、回环验证及地图是否同步更新 \\

轨迹整体可任意移动 & 全局参考或先验没有固定 \\

\bottomrule\end{tabular}\end{center}

这些检查把第六章的局部微分、第七章的概率、第八九章的状态误差，以及第十十一章的相对约束放在同一处。一个数值结果只有连同它的坐标、时间、线性化点和不确定性一起保存，后续模块才知道该怎样使用它。

\section{最终记忆框架}

\begin{equation*}
\begin{aligned}
\text{IMU与LiDAR原始数据}
&\longrightarrow\text{前端当前状态、去畸变和局部配准}\\
&\longrightarrow\text{关键帧、IMU预积分与几何约束}\\
&\longrightarrow\text{后端联合优化历史状态}\\
&\longrightarrow\text{回环检索、验证与全局校正}\\
&\longrightarrow\text{更新地图和实时参考}.
\end{aligned}
\end{equation*}

前端求及时，后端求多份历史信息之间的一致，回环补上长时间之后重新看到同一地点的约束。三者共享同一个基本做法：先弄清状态和观测分别是什么，在合法的局部坐标中构造残差与微分，再根据噪声模型决定每份数据对估计有多大影响。
